\documentclass[10pt]{article}
\usepackage[margin=20mm]{geometry}
\usepackage{amsmath,amssymb,mathtools,fontspec,unicode-math,xurl,hyperref}
\IfFontExistsTF{DejaVu Serif}{\setmainfont{DejaVu Serif}}{\setmainfont{Latin Modern Roman}}
\setmathfont{latinmodern-math.otf}
\setlength{\parindent}{0pt}\setlength{\parskip}{6pt}\sloppy
\title{Technical Review of
Planar Regularity Structures}
\author{AI-assisted technical review}
\date{}
\begin{document}
\maketitle
Abstract: This report checks the grading and algebraic model construction in Section 3 of Ludwig Rahm’s Planar Regularity Structures. Two explicit admissible trees show that Lemma 3.5, with the definitions actually printed, does not establish the claimed regularity structure: one assigned degree is absent from the stated index set, and an allowed resonant regularity produces a nonconstant degree-zero basis element. The algebraic model identities are examined separately and are not refuted by these examples. Both findings concern a central construction; their supported assessment is Major. The later deformed Hopf-algebra theory is not certified or refuted here. Scientific adjudication remains pending.\par
\section*{Identity, version, and inspected scope}
The reviewed work is Ludwig Rahm, Planar Regularity Structures, Foundations of Computational Mathematics 26 (2026), 1269–1319, DOI https://doi.org/10.1007/s10208-025-09701-1 [rahm]. This is the 51-page publisher version of record, published online 8 September 2025 and acquired on 27 October 2026. Its exact bytes, rather than the earlier 37-page arXiv version, define this review. Here PDF page means a one-based physical page; printed page \(1268+j\) corresponds to PDF page \(j\).\par
The focused checks cover Definitions 2.8–2.9, the definition of \(\mathcal{PB}\) and its degree in Section 3, Lemmas 3.3 and 3.5, Proposition 3.6, and the algebraic identities in Theorem 3.8. The introductory account, the setup of Section 4, its discussion of restricted tree spaces, and its final extended-decoration argument were read to bound the impact. The long proofs of the deformed positive and negative Hopf structures and the complete cointeraction calculation were not independently verified. Source text was parsed on all 51 pages; that operation is not claimed as scientific reading of all pages. The coverage ledger records the actual page ranges and unexamined material.\par
\section*{Definitions and assumptions relevant to the findings}
On PDF pages 16–17 (printed 1284–1285), Definition 2.8 requires \[A\subset\mathbb R,\qquad 0\in A,\qquad
 T=\bigoplus_{\beta\in A}T_\beta,\qquad T_0\simeq\mathbb R,\] with a lower-bounded locally finite index set, Banach homogeneous spaces, and a triangular structure-group action. Definition 2.9 explicitly sets aside the usual analytic model bounds. Consequently, the absence of such bounds is not treated as a defect.\par
Equation reference: eq:printedA. On PDF page 17, \(\mathcal{PB}\) is the vector space of planar rooted trees with edge labels \(0,1,\ldots,n\). A zero label is represented by an undecorated edge. Each nonzero edge ends at a leaf; a vertex has at most one such decorated edge, and that edge is rightmost. On PDF page 18 (printed 1286), the specified index set is \[
 A=\{0\}\cup_i\alpha_i\mathbb N
       \cup_i(\alpha_i\mathbb N-1),\] and the degree of a tree is the sum of \(1\) for every undecorated edge and \(\alpha_i-1\) for every edge decorated by \(i\). The formula is visible in the original PDF; it is not an extraction substitution for an additive sum of sets. No reciprocal-integer or nonresonance restriction on the positive regularities is imposed in this construction. The preceding rough-path definition permits regularity parameters in \((0,1]\) (PDF pages 10–11).\par
Lemma 3.5, on PDF page 20, asserts that \((A,\mathcal{PB},G)\) is a regularity structure. Its proof assumes that the indicated degree makes \(\mathcal{PB}\) an \(A\)-graded space, then checks triangularity. The following counterexamples test this algebraic assertion. They require no stochastic sample, numerical solver, or extra assumption about a particular rough-path lift. In particular, the time-augmentation and integration condition (17), introduced for the later model, cannot remove basis trees from the already declared space in Lemma 3.5.\par
\section*{Substantive findings}
Problem 1 (The printed index set omits an allowed tree degree).  For one nonzero edge label, set \(\alpha_1=3/5\). Let \(\tau\) be the chain with four vertices \(v_0,v_1,v_2,v_3\), root \(v_0\), and successive edge labels \(0,0,1\). Then \(\tau\in\mathcal{PB}\) but the assigned degree of \(\tau\) does not belong to [eq:printedA] (\#eq:printedA). Thus the first grading assertion in the proof of Lemma 3.5 is false for this allowed parameter.\par
\subsection*{Complete check}
Every nonroot vertex of the chain has exactly one parent. The only nonzero edge ends at the leaf \(v_3\); it is the only outgoing edge of \(v_2\), and hence is rightmost. There is at most one decorated edge incident to any vertex. These are exactly the tree restrictions on PDF page 17. The chain has positive degree and no decorated edge at the root, so restricting attention to the positive tree space would not remove this example.\par
Equation reference: eq:degree1. The displayed degree rule gives \[
 |\tau|=1+1+\left(\frac35-1\right)=\frac85.\] For \(8/5\) to lie in either family supplied by \(\alpha_1\), an integer \(k\in\mathbb N\) would have to satisfy \[\frac35k=\frac85 \quad\hbox{or}\quad
 \frac35k-1=\frac85.\] The two required values are \(k=8/3\) and \(k=13/3\), neither an integer. Including a time parameter \(\alpha_0=1\), the interpretation most favorable to the assertion, adds only integer degrees, which still do not include \(8/5\). The argument is unaffected by whether \(\mathbb N\) includes zero. Since \(\tau\) is a nonzero tree basis element with the degree assigned in [eq:degree1] (\#eq:degree1), it cannot be placed in any homogeneous component indexed by [eq:printedA] (\#eq:printedA) without changing the printed grading.\par
The retained independent Python check uses exact rational arithmetic, verifies the rooted-tree, leaf, and planar-order conditions, and solves these membership equations. It supplies a finite verification of this witness, while the preceding calculation is the complete argument.\par
\subsection*{Alternative interpretations and impact}
If all nonzero regularities are chosen as one reciprocal integer \(1/N\), the particular missing-degree obstruction disappears because adding an integer preserves that lattice. That is a valid special case, not an assumption stated in Lemma 3.5. Reading the unions as the additive semigroup generated by the regularities would also change the printed definition. The source image and publisher HTML both show unions.\par
This is a failure of the specified regularity-structure construction: the model theorem explicitly uses that triple. The supported assessment is Major, because Lemma 3.5 supplies a central structural prerequisite for Theorem 3.8. The example does not show that the recentering identities are false, that a corrected grading cannot work, or that the later Hopf-algebra identities fail.\par
Problem 2 (A permitted resonant regularity gives a nontrivial zero sector).  Set \(\alpha_1=1/2\). Consider the planar tree \(\sigma\) with edges \[(v_0,v_1;0),\qquad (v_0,v_2;1),\qquad (v_1,v_3;1),\] where \(v_0\) is the root and the edge to \(v_2\) is to the right of the edge to \(v_1\). Then \(\sigma\) is a nonconstant degree-zero basis element of \(\mathcal{PB}\). Together with the unit tree, it contradicts \(T_0\simeq\mathbb R\) in the paper’s Definition 2.8.\par
\subsection*{Complete check}
Equation reference: eq:degree2. The vertices \(v_2,v_3\) are leaves and receive the two decorated edges. Each of \(v_0,v_1\) has one decorated outgoing edge. At \(v_0\) it is rightmost by construction, and at \(v_1\) it is the sole outgoing edge. Every other restriction of \(\mathcal{PB}\) is satisfied. The degree rule gives \[
 |\sigma|=1+\left(\frac12-1\right)+
                 \left(\frac12-1\right)=0.\] The one-vertex unit tree also has degree zero. The four-vertex tree \(\sigma\) and the one-vertex unit are distinct basis elements in the vector space defined on PDF page 17. No relation identifying them, or quotient deleting \(\sigma\), is imposed there. Hence \[\dim\bigl(\mathcal{PB}_0\bigr)\ge 2,\] whereas the stated Definition 2.8 requires dimension one. No simulation or claimed analytic model evaluation is needed for this dimension argument.\par
\subsection*{Alternative interpretations and impact}
A convention in which the zero sector merely contains a distinguished copy of \(\mathbb R\) would avoid this specific contradiction, but it is different from the definition used by this paper. Alternatively, carefully chosen nonresonant degrees might exclude nonunit zero-degree trees. Such a hypothesis is absent from Lemma 3.5. A later equation-specific restriction is not specified for the full Section 3 space. This remains a separate obstruction even after enlarging the index set.\par
The supported assessment is again Major, with the same bounded dependency: Lemma 3.5 does not provide a regularity structure under its own definition for all allowed parameters. This finding does not establish failure of the model’s algebraic identities or of an appropriately restricted construction.\par
\section*{Checks that survive and the precise repair boundary}
\subsection*{The model identities}
Theorem 3.8 assumes time augmentation and condition (17) (PDF pages 20–22). On an element \(B_+\varphi(\omega)\), the proof reduces the model compatibility equation to \[\langle X_{yx}\otimes X_{xt},\Delta_{\mathrm{MKW}}\omega\rangle
 =\langle X_{yx}*X_{xt},\omega\rangle
 =\langle X_{yt},\omega\rangle,\] which is Chen’s identity with the indicated order of endpoints. For a tree with a decorated root edge, the differentiated constant term vanishes because Lemma 3.7 gives \(\Pi_x(1)=1\). Condition (17) is needed for the derivative identification and is explicitly assumed. These steps were read and checked; the two grading counterexamples do not refute them. They prevent those identities, as written, from completing the claimed model on the uncorrected regularity-structure triple.\par
\subsection*{What an index-set change actually restores}
There is a precise limited correction to Problem 1: replace [eq:printedA] (\#eq:printedA) by the set of degrees of the declared admissible trees, including the unit. For finitely many strictly positive \(\alpha_i\), this degree set is lower bounded and locally finite.\par
Indeed, write \(k\) for the number of undecorated edges and \(m_i\) for the number of edges labeled \(i\), with \(m=\sum_i m_i\). There are \(k+1\) non-noise vertices. The one-noise-edge-per-vertex condition implies \(m\le k+1\). Therefore \[|\eta|=k-m+\sum_i m_i\alpha_i,\qquad k-m\ge-1.\] For \(|\eta|\le R\), the nonnegative counts obey \(\sum_i m_i\alpha_i\le R+1\), so there are only finitely many choices of the \(m_i\), and then only finitely many choices of \(k-m\). The nonunit degrees are bounded below by \(\min(0,\min_i\alpha_i-1)>-1\). This proves the stated index-set properties.\par
This change restores degree coverage and these two index-set properties only. It does not cure Problem 2, choose a permissible quotient, establish all model-space axioms, or constitute a verified replacement for the full Lemma 3.5. Repair of the zero-sector condition and its compatibility with the coaction remains open in this review.\par
\subsection*{Later construction and versions}
Section 4 defines its index set as the actual set of regularities (PDF page 26), and explicitly discusses restriction to suitable tree subspaces (PDF page 36). The extended-decoration argument in Proposition 4.28 preserves the degree used in positive projection (PDF page 49). These are reasons to confine the present findings to the stated Section 3 construction rather than infer a refutation of all subsequent algebra.\par
A bounded check of the publisher page and the arXiv record on 29 October 2026 found no linked correction addressing these points. The arXiv record listed only the earlier version v1 [arxiv]. Its full mathematical content was not compared with the acquired publisher PDF. This bounded search does not establish first discovery or the absence of unpublished corrections. No author contact or author-code execution occurred.\par
\section*{Conclusion and distinct states}
Item | Issue assessment | Repair and remaining work\par
Problem 1 | Major; missing allowed degree | Actual degree set repairs degree coverage; full construction not re-proved.\par
Problem 2 | Major; zero-sector dimension | A compatible restriction or revised definition remains to be supplied.\par
Model identities | No issue established in the checked algebraic steps | Retain condition (17) and time augmentation; this is not a complete analytic model verification.\par
Scientific decision: pending. Human reviews: none supplied. Author adoption: unknown and not requested. This is an initial review; independent acceptance has not been performed.\par
\section*{Supporting checks}
The exact-rational checker is `grading\_check.py`. It is independently written; no author program was executed. The admissibility and degree calculations are analytical arguments in this report supported by that finite check.\par
Scientific reading, argument assessment, report preparation, compilation, parsing and actual page inspection were performed in the assigned execution conversation. The main dispatcher’s product-interface observation is attributed to that dispatcher; the execution environment’s self-description is recorded separately in provenance without a claim of backend equivalence. An internal helper supplied only a clerical inventory checklist. There is no human review attestation.\par
9 L. Rahm, Planar Regularity Structures, Found. Comput. Math. 26 (2026), 1269–1319. Publisher version of record, online 8 September 2025. https://doi.org/10.1007/s10208-025-09701-1. L. Rahm, Planar Regularity Structures, arXiv:2212.04856, version-history metadata consulted 29 October 2026; v1 submitted 9 December 2022. Only the record was consulted for version history. https://arxiv.org/abs/2212.04856.\par
\end{document}
